
\begin{equation}\tag{1.1.3}
\cl_\bullet(\underline{F})=\sum_{x\in X(n)}\longg(\underline{F}_x)\cl_\bullet(\underline{O}_{Z(x)})+w
\end{equation}
\uline{où} $w\in\Filt_{n-1}(X)$.

On en déduit par récurrence immédiate sur $n$ :

\textbf{Corollaire 1.1.4.} \uline{Pour tout schéma noethérien} $X$ \uline{et pour tout entier} $n$, $\Filt_n(X)$ \uline{est le sous-groupe de} $K_\bullet(X)$ \uline{engendré par les classes} $\cl_\bullet(\underline{O}_Z)$\uline{, où} $Z$ \uline{parcourt l'ensemble des sous-schémas intègres fermés de dimension} $\leq n$ \uline{de} $X$.

\emph{Démonstration de} 1.1.2 : Soit $X'=\Supp(\underline{F})$. Comme $X'$ est un espace noethérien de dimension $\leq n$, les éléments de $X'\cap X(n)$ sont maximaux dans $X'$, donc en nombre fini, et $\longg(\underline{F}_x)$ est finie pour tout $x\in X'\cap X(n)$. Cela prouve la première assertion.

Soit ${K}=\Coh(X)$ la catégorie des $\underline{O}_X$-modules cohérents et soit ${K}'$ le sous-ensemble de $\operatorname{Ob}({K})$ formé des $\underline{O}_X$-modules cohérents concentrés sur $X'$ qui satisfont à la condition (1.1.3). Évidemment ${K}'$ contient les faisceaux $\underline{O}_{Z(x)}$ avec $x\in X'$. Il suffit donc de montrer que ${K}'$ est un sous-ensemble «~exact~» de $\operatorname{Ob}({K})$, car alors le lemme de dévissage (EGA III 3.1.2) donne $\underline{F}\in{K}'$.

Soit donc
\begin{equation}\tag{1.1.5}
0\longrightarrow \underline{G}'\longrightarrow G\longrightarrow \underline{G}''\longrightarrow0
\end{equation}
une suite exacte de $\underline{O}_X$-modules ayant deux des membres dans ${K}'$. Alors les faisceaux $\underline{G}'$, $\underline{G}$ et $\underline{G}''$ sont tous concentrés sur $X'$, donc les longueurs $\ell'_x$, $\ell_x$ et $\ell''_x$ des $\underline{O}_{X,x}$-modules $\underline{G}'_x$, $\underline{G}_x$ et $\underline{G}''_x$ sont finies
